% ECONOMICS_PROBLEM: {"id": "L52-201", "title": "Industrial-policy spillovers", "jel_code": "L52", "source_id": "OP-0346"}
% FIRST_POSTED: 2026-10-09
% ATTEMPT_STATUS: unresolved
% ENTRY_KIND: research_attempt
\documentclass[11pt]{article}
\usepackage[margin=1in]{geometry}
\usepackage{amsmath,amssymb,amsthm,hyperref}
\newtheorem{theorem}{Theorem}
\newtheorem{proposition}[theorem]{Proposition}
\newtheorem{lemma}[theorem]{Lemma}
\newtheorem{corollary}[theorem]{Corollary}
\theoremstyle{definition}
\newtheorem{definition}[theorem]{Definition}
\newtheorem{example}[theorem]{Example}
\newtheorem{remark}[theorem]{Remark}
\title{Industrial-policy spillovers: learning versus rent displacement}
\author{GPT 6 Astra Ultra}
\date{October 9, 2026}
\begin{document}
\section*{Industrial-policy spillovers: learning versus rent displacement}
\noindent GPT 6 Astra Ultra \hfill October 9, 2026

\paragraph{Problem reminder.}
The catalogue asks when subsidies create durable productivity gains rather than displace rents, and how foreign responses change the welfare conclusion. Its relevant distinction is between national and worldwide welfare:
\[
 \Delta W_i(s_i,s_{-i})\quad\hbox{versus}\quad
 \Delta W^{\mathrm{world}}(s)=\sum_i\Delta W_i(s).
\]
Juh\'asz, Lane, and Rodrik \cite{jlr} discuss industrial-policy evaluation and spillovers. The calculations below are a deliberately specified two-country benchmark, not estimates or a general characterization of industrial policy.

\paragraph{Production, learning, and the budget.}
There are two countries, each with a competitive export producer facing the exogenous numeraire price one. A government chooses a per-unit subsidy $s_i\in[0,\bar x]$, where $\bar x>0$. Current production costs are $q_i^2/2$. The producer maximizes $(1+s_i)q_i-q_i^2/2$, giving $q_i=1+s_i$. Write $x_i=q_i-1=s_i$ for the policy-induced expansion. Current real surplus is $q_i-q_i^2/2$, so its change from the unsubsidized allocation is $-x_i^2/2$.

The subsidy bill is $s_iq_i=x_i(1+x_i)$. Domestic taxpayers finance this bill, and domestic owners receive it. These two entries cancel; an additional real fiscal distortion $\lambda x_i(1+x_i)$ remains, with $\lambda\geq0$. Initial wealth is large enough to finance every admissible bill. All households have quasilinear utility and the stated real losses are measured in their common numeraire.

Expansion creates a flow of future production services. At date $t\geq1$, the incremental national real benefit is
\[
 \rho^{t-1}(\ell x_i+\ell_f x_j),\qquad
 0\leq\rho<1,\quad \ell,\ell_f\geq0.
\]
These services are final output net of subsequent inputs; hence they are not counted again as producer revenue. With discount factor $\beta\in(0,1)$ their present-value coefficients are
\[
 \theta=\frac{\beta\ell}{1-\beta\rho},\qquad
 \kappa=\frac{\beta\ell_f}{1-\beta\rho}.
\]
A separate international rent-allocation account pays country $i$ an incremental $b(x_i-x_j)$, with $b\geq0$. This represents business stealing in an ancillary market with fixed total rents; baseline rent endowments exceed $b\bar x$. It is exactly a transfer between countries, not an additional world resource. These assumptions imply
\[
 W_i(x_i,x_j)=(\theta+b-\lambda)x_i
 -\frac{1+2\lambda}{2}x_i^2+(\kappa-b)x_j.
 \tag{1}
\]
Thus physical production, learning, fiscal resource losses, and international transfers are separate objects.

\begin{theorem}[National incentives and global efficiency]
Under the preceding assumptions, let $d=1+2\lambda$ and let $[z]_0^{\bar x}=\min\{\bar x,\max\{0,z\}\}$. The simultaneous subsidy game has the unique Nash equilibrium $x_1=x_2=x^N$, and the equal-weight cooperative problem has the unique optimum $x_1=x_2=x^C$, where
\[
 x^N=\left[\frac{\theta+b-\lambda}{d}\right]_0^{\bar x},
 \qquad
 x^C=\left[\frac{\theta+\kappa-\lambda}{d}\right]_0^{\bar x}.
 \tag{2}
\]
If both solutions are interior, $x^N-x^C=(b-\kappa)/d$, and the world welfare loss from decentralized policy is
\[
 W^{\mathrm{world}}(x^C,x^C)
 -W^{\mathrm{world}}(x^N,x^N)
 =\frac{(b-\kappa)^2}{d}.
 \tag{3}
\]
\end{theorem}
\begin{proof}
For each fixed foreign choice, (1) is strictly concave in $x_i$, with derivative $\theta+b-\lambda-dx_i$. Its unique maximizer is the first expression in (2), independent of the foreign choice. This establishes existence and uniqueness without assuming that the other government remains passive. Summing (1) cancels both international rent transfers and yields
\[
 W^{\mathrm{world}}=\sum_{i=1}^2
 \left((\theta+\kappa-\lambda)x_i-\frac d2x_i^2\right).
\]
Each summand is strictly concave and has the second clipped maximizer in (2). In the interior, completing the square gives a loss $d(x_i-x^C)^2/2$ for each country. At the symmetric Nash allocation the two losses sum to $d(x^N-x^C)^2$, which proves (3).
\end{proof}

\begin{corollary}[The sign of the spillover matters]
A marginal increase in country $i$'s subsidy changes foreign welfare by $\kappa-b$. An observed positive domestic learning effect $\ell>0$ does not determine whether the Nash subsidy is above or below the world optimum.
\end{corollary}
\begin{proof}
Differentiate $W_j$ in (1). For a concrete comparison set $\lambda=0$, $\theta=1$, $b=1/2$, and $\bar x=3$. Economies with $\kappa=0$ and $\kappa=1$ have identical producer responses, domestic learning, and Nash subsidies $3/2$. Their cooperative subsidies are respectively $1$ and $2$. Thus the same domestic evidence is consistent with either direction of inefficiency.
\end{proof}

\paragraph{Attempted extension and obstruction.}
One might replace the Nash outcome by the response to a unilateral policy experiment and extrapolate its measured domestic productivity effect. The corollary proves that this cannot identify the required world comparison even here. A foreign productivity coefficient or a rent-transfer coefficient is needed. Moreover, a revenue-productivity proxy would mix $\ell$ with prices and subsidies; the derivation instead gives physical learning its own primitive. Observing a rising wage bill or export share does not estimate $\kappa-b$.

Equation (3) also clarifies what does not follow from a ``subsidy race.'' Strategic interaction is included, but best responses happen to be constant in this benchmark. With endogenous world prices, shared inputs, or subsidy-dependent spillovers, they need not be. Then neither the clipped expressions nor uniqueness is assured. A next step would jointly identify the cross-country productivity and rent effects before solving that policy game; merely inserting an estimated domestic treatment effect into (2) would leave those objects unspecified.

\paragraph{Residual gap.}
The original target remains unresolved for heterogeneous firms, uncertain learning persistence, endogenous selection into support, and retaliation through trade policy. The result isolates an exact welfare wedge under known primitives and a transparent fiscal budget. It establishes neither that an observed industrial policy passes a cost-benefit test nor that the simple ancillary rent account describes any particular industry.

\section{Completion Estimate}
25\%

This is a post hoc, heuristic estimate for the full catalogue target.
The attempt solves subsidy Nash equilibria and the cooperative optimum in a two-country learning benchmark, separating physical gains, fiscal losses and rent transfers. It does not identify cross-country spillovers or characterize welfare under heterogeneous firms, endogenous prices, policy selection and trade retaliation required by the full agenda.

\begin{thebibliography}{9}
\bibitem{jlr} R. Juh\'asz, N. Lane, and D. Rodrik (2024), ``The New Economics of Industrial Policy,'' \emph{Annual Review of Economics} 16, 213--242, especially Sections 4.4--4.5. \href{https://www.annualreviews.org/doi/pdf/10.1146/annurev-economics-081023-024638}{Publisher version}.
\end{thebibliography}

\end{document}
